% Active Inference Pong — Mathematics
% Converted from PONG.md
\documentclass[12pt, a4paper]{article}

% ── Packages ──────────────────────────────────────────────────────────────────
\usepackage[utf8]{inputenc}
\usepackage[T1]{fontenc}
\usepackage{amsmath, amssymb, amsthm}
\usepackage{bm}
\usepackage{booktabs}
\usepackage[margin=2.5cm]{geometry}
\usepackage[colorlinks=true, linkcolor=blue!60!black, urlcolor=blue!60!black]{hyperref}
\usepackage{xcolor}
\usepackage{listings}
\usepackage{parskip}
% microtype removed (font expansion unsupported on this TeX install)
\usepackage{pifont}   % for \ding{51} checkmark
\usepackage{array}

% ── Custom commands ───────────────────────────────────────────────────────────
\newcommand{\E}{\mathbb{E}}
\newcommand{\R}{\mathbb{R}}
\newcommand{\1}{\mathbf{1}}
\newcommand{\cN}{\mathcal{N}}
\newcommand{\cW}{\mathcal{W}}
\newcommand{\cNIW}{\mathcal{NIW}}
\newcommand{\Dir}{\mathrm{Dir}}
\newcommand{\tr}{\operatorname{tr}}
\newcommand{\clip}{\operatorname{clip}}
\newcommand{\sgn}{\operatorname{sgn}}
\newcommand{\sect}[1]{\S\ref{#1}}  % section cross-reference helper
\newcommand{\code}[1]{\texttt{#1}}
\newcommand{\yes}{\ding{51}}

% ── Listings style for code blocks ───────────────────────────────────────────
\lstset{
  basicstyle=\small\ttfamily,
  breaklines=true,
  frame=single,
  backgroundcolor=\color{gray!8},
  framerule=0pt,
}

% ── Section numbering starting from 0 ────────────────────────────────────────
\setcounter{section}{-1}

% ── Title ─────────────────────────────────────────────────────────────────────
\title{\textbf{Active Inference Pong --- Mathematics}\\[0.5em]
  \large A gradient-free, CAVI-based Active Inference agent for ALE/Pong-v5}
\author{}
\date{}

% =============================================================================
\begin{document}
% =============================================================================

\maketitle

\begin{center}
\fbox{\parbox{0.92\textwidth}{%
  \textbf{Current implementation:} MPPI receding-horizon planner with greedy
  nominal, adaptive EMA target smoothing, urgency-scaled preference covariance,
  opponent-aware placement, and cross-session Bayesian aim calibration.
}}
\end{center}

\tableofcontents
\newpage

% =============================================================================
\section{Problem and architecture}
\label{sec:0}
% =============================================================================

\subsection{Problem description}
\label{sec:0.1}

\textbf{ALE/Pong-v5} is a two-player video game in which each player controls a
vertical paddle.  The right player (agent) must intercept a moving ball with its
paddle and return it past the left opponent's paddle to score.  The first to reach
21 points wins; an episode ends when either player reaches that score.

The agent reads four bytes directly from the Atari~2600 RAM each frame:

\begin{table}[ht]
\centering
\begin{tabular}{lcc}
\toprule
Signal & RAM byte & Range \\
\midrule
Ball $x$ & 49 & 0--255 \\
Ball $y$ & 54 & 0--255 \\
Player paddle $y$ & 51 & 0--255 \\
Opponent paddle $y$ & 50 & 0--255 \\
\bottomrule
\end{tabular}
\end{table}

All values are normalised to $[0,1]$.  The \textbf{ball velocity} is not available
in any single frame --- it is a latent variable that must be inferred from the
sequence of observations.

\textbf{What makes this hard:}
\begin{itemize}
  \item \textit{Latent dynamics.}  Neither velocity nor acceleration is directly
    observable.  The agent must maintain a probabilistic belief over the full 6D
    state $[b_x, b_y, \dot b_x, \dot b_y, p_y, o_y]$ and propagate it forward in
    time.
  \item \textit{Non-linear bounce dynamics.}  The ball reflects off the top and
    bottom walls (sign flip in $\dot b_y$) and off both paddles (sign flip in
    $\dot b_x$).  A single Gaussian cannot represent the bimodal predictive
    distribution at bounce events.
  \item \textit{Adversarial opponent.}  The opponent tracks the ball reactively.
    Winning requires not just intercepting the ball but placing the return where
    the opponent cannot reach it.
  \item \textit{Tight timing.}  The paddle moves $\approx 0.060$ normalised units
    per frame (frame-skip~$=4$).  Missing a rally by one paddle-step is enough to
    lose the point.
\end{itemize}

\textbf{What the agent does NOT use:}
\begin{itemize}
  \item No reward signal, no policy network, no value function.
  \item No replay buffer, no gradient descent, no neural networks.
  \item No hand-crafted game-specific heuristics beyond the choice of RAM bytes.
\end{itemize}

\textbf{What the agent DOES use:}

The agent is built on \textbf{Active Inference} (Friston et al.) --- a framework
in which behaviour emerges from minimising a single objective, the
\textbf{Expected Free Energy} (EFE).  At each frame the agent asks:
\textit{which action leads to a future that best matches my preferred
observations?}  Preferences encode the goal (intercept the ball, put it past the
opponent); the world model (VBGS) encodes how the game evolves; MPPI searches for
the action sequence that minimises EFE over a 15-frame horizon.

\textbf{Current performance} (frozen VBGS, rMM win-rate model, 30 episodes):

\begin{table}[ht]
\centering
\begin{tabular}{lc}
\toprule
Agent & Mean score \\
\midrule
Rule-based (track ball $y$) & $-14.4$ \\
EFE $H=3$, intercept target & $-6.6$ \\
EFE frozen, rMM, 30 episodes & $-7.4$ \\
\bottomrule
\end{tabular}
\end{table}

The EFE agent beats the rule-based baseline by \textbf{$+7$ points} on average
and regularly wins individual episodes.

% ─────────────────────────────────────────────────────────────────────────────
\subsection{Block diagram}
\label{sec:0.2}
% ─────────────────────────────────────────────────────────────────────────────

\begin{figure}[ht]
\centering
\begin{lstlisting}[basicstyle=\scriptsize\ttfamily, frame=none, backgroundcolor=\color{white}]
+------------------------------------------------------+
|                 ALE / Pong-v5                        |
|           Atari RAM, frame-skip = 4                  |
+----------------------+-------------------------------+
                       | o_t = [bx,by,py,oy]/255
                       v
      +----------------------------------------+
      |          World model                   |
      |  BallTransitionVBGS  K=3 NIW ........  |--[online CAVI]--+
      |  OpponentTransVBGS   K=2 NIW           |                 |
      |    +-> OpponentBeliefTracker            |--[innovations]--+---> R
      |  LikelihoodModel:  o = H*s + eps        |
      |  MixtureBeliefFilter  K=2 Kalman        |
      +--------------------+-------------------+
                           | mu_t, Sigma_t (solid)
                           v
      +----------------------------------------+
      |     PriorPreferences                   |<--- AimCalibrator
      |  Landing predictor -> y_land           |       (6 EMAs)
      |  Urgency scaling of sigma_py           |
      |  Opponent-aware placement: y* += delta |<--- Win-rate rMM
      +--------------------+-------------------+     (NIW+Beta, K<=16)
                           | o*, C^{-1} (solid)
                           v
      +----------------------------------------+
      |      MPPI Planner  (EFE Agent)         |
      |  H=15 frames,  N_r=128 rollouts        |
      |  VBGS step-0 + linear steps 1..H-1     |
      |  soft-min over rollout EFE costs       |
      +--------------------+-------------------+
                           | a_t in {0=NOOP, 2=UP, 3=DOWN}
                           v
+------------------------------------------------------+
|                 ALE / Pong-v5                        |
|   (rally outcomes --[dashed]--> rMM & AimCalibrator) |
+------------------------------------------------------+
\end{lstlisting}
\caption{Architecture of the Active Inference Pong agent. Solid arrows: signal
  flow each frame. Dashed connections: online learning flow.}
\label{fig:architecture}
\end{figure}

\textbf{Signal flow (solid arrows):} each frame, the environment emits a RAM
observation $\to$ the Belief Filter predicts and corrects to produce
$(\mu_t, \Sigma_t)$ $\to$ the MPPI planner evaluates 128 perturbed rollouts
against the preference target $\to$ the lowest-EFE first action is sent back to
the environment.

\textbf{Learning flow (dashed arrows):} the VBGS and likelihood noise $R$ are
refined online from live transitions and Kalman innovations; the win-rate rMM
and aim calibrator accumulate rally outcomes across sessions.

% =============================================================================
\section{State and observations}
\label{sec:1}
% =============================================================================

The agent reads a small number of bytes from the Atari~2600 RAM on every simulated
frame, rather than processing raw pixel images.  This section defines those signals
and the 6D latent state built on top of them.  The key design constraint is that
\textbf{velocities are not available} from a single frame: they must be inferred by
the belief filter (\S\ref{sec:9}) from the sequence of positions.  This is why the
state is strictly larger than the observation --- the observation tells us where
things are, but the state also encodes how fast they are moving.

All RAM bytes are normalised to $[0,1]$ by dividing by 255.

\subsection{Ball state}
\[
s^b_t = \begin{bmatrix} b_x \\ b_y \\ \dot b_x \\ \dot b_y \end{bmatrix} \in \R^4
\]
Velocities are computed by frame differencing:
\[
\dot b_t = b_t - b_{t-1}
\]

\subsection{Opponent state}
\[
s^o_t = \begin{bmatrix} o_y \end{bmatrix} \in \R^1
\]

\subsection{RAM addresses used}

\begin{table}[ht]
\centering
\begin{tabular}{lcc}
\toprule
Symbol & RAM index & Meaning \\
\midrule
$b_x$ & 49 & Ball $x$ position \\
$b_y$ & 54 & Ball $y$ position \\
$p_y$ & 51 & Player paddle $y$ (right) \\
$o_y$ & 50 & Opponent paddle $y$ (left) \\
\bottomrule
\end{tabular}
\end{table}

The four observable positions combine with two latent velocities to form the
\textbf{full 6D state}
$s = [b_x,\, b_y,\, \dot b_x,\, \dot b_y,\, p_y,\, o_y]^\top$ used throughout
\S\ref{sec:8}--\S\ref{sec:9}.  The transition models (\S\ref{sec:5}) learn to
predict $s_{t+1}$ from $s_t$; the likelihood model (\S\ref{sec:8}) connects the
6D state back to the 4D observation via the $H$ matrix.

% =============================================================================
\section{Normal-Inverse-Wishart prior}
\label{sec:2}
% =============================================================================

Every component of the mixture models in this system --- the VBGS transition model
(\S\ref{sec:3}), the win-rate rMM (\S\ref{sec:17}) --- stores its sufficient
statistics as a Normal-Inverse-Wishart (NIW) distribution over the component's
mean and precision matrix.  The NIW family is the \textbf{conjugate prior} to the
multivariate Gaussian likelihood: after observing data, the posterior has the same
parametric form as the prior, and the update reduces to simple arithmetic on four
numbers $(\mathbf{m}_k, \kappa_k, \nu_k, \mathbf{W}_k)$ rather than an integral.
Conjugacy is the reason the whole system is gradient-free --- every learning step
is a closed-form update, not a gradient-descent step.

Each mixture component $k$ has a Gaussian likelihood with unknown mean $\mu_k$ and
precision $\Lambda_k$.  The conjugate prior is:
\[
p(\mu_k, \Lambda_k) = \cNIW(\mathbf{m}_k,\, \kappa_k,\, \nu_k,\, \mathbf{W}_k)
\]

\begin{table}[ht]
\centering
\begin{tabular}{ll}
\toprule
Symbol & Meaning \\
\midrule
$\mathbf{m}_k \in \R^d$ & prior mean \\
$\kappa_k > 0$           & mean pseudo-count \\
$\nu_k > d - 1$          & degrees of freedom \\
$\mathbf{W}_k \in \R^{d\times d}$ & scale matrix \\
\bottomrule
\end{tabular}
\end{table}

\subsection{Variational expectations}

Under the variational posterior $q(\Lambda_k) = \cW(\nu_k, \mathbf{W}_k)$:
\[
\E\!\left[\log|\Lambda_k|\right]
= \sum_{i=0}^{d-1} \psi\!\left(\frac{\nu_k - i}{2}\right)
  + d\ln 2 + \ln|\mathbf{W}_k|
\]
where $\psi$ is the digamma function.  Under
$q(\mu_k, \Lambda_k) = \cNIW(\mathbf{m}_k, \kappa_k, \nu_k, \mathbf{W}_k)$:
\[
\E\!\left[(\mathbf{x} - \mu_k)^\top \Lambda_k\,(\mathbf{x} - \mu_k)\right]
= \frac{d}{\kappa_k}
  + \nu_k\,(\mathbf{x} - \mathbf{m}_k)^\top
    \mathbf{W}_k^{-1}(\mathbf{x} - \mathbf{m}_k)
\]

\subsection{Posterior update (CAVI M-step)}

Given effective count $N_k$, weighted mean $\bar{\mathbf{x}}_k$, and scatter
$\mathbf{S}_k$:
\[
\kappa_k^* = \kappa_k + N_k
\qquad
\mathbf{m}_k^* = \frac{\kappa_k \mathbf{m}_k + N_k \bar{\mathbf{x}}_k}{\kappa_k^*}
\]
\[
\nu_k^* = \nu_k + N_k
\]
\[
\mathbf{W}_k^* = \mathbf{W}_k + \mathbf{S}_k
+ \frac{\kappa_k N_k}{\kappa_k^*}\,
  (\bar{\mathbf{x}}_k - \mathbf{m}_k)(\bar{\mathbf{x}}_k - \mathbf{m}_k)^\top
\]

\subsection{Posterior predictive}

\textbf{Two stacked approximations apply here --- neither is exact.}

\textbf{Approximation 1 --- variational M-step (CAVI).}  In a mixture model the
true posterior over parameters is not NIW because the latent assignments
$\mathbf{Z}$ couple all components.  CAVI imposes the mean-field factorisation
$q = q(\mathbf{Z})\prod_k q(\mu_k,\Lambda_k)$.  The conjugate prior guarantees
the \textit{form} of each factor remains NIW and the M-step stays closed-form, but
$\mathbf{W}_k^*$ is the scale of the \textbf{variational} posterior, not the exact
Bayesian posterior.

\textbf{Approximation 2 --- Gaussian predictive.}  The exact posterior predictive
under the NIW is a multivariate Student-$t$ with scale:
\[
\Sigma_\text{pred}^{\text{exact}}
= \frac{\kappa_k+1}{\kappa_k(\nu_k - d + 1)}\,\mathbf{W}_k
\]
This implementation uses the Inverse-Wishart posterior mean instead (a Gaussian
approximation):
\[
\hat{\Sigma}_k = \frac{\mathbf{W}_k}{\nu_k - d - 1} \qquad (\nu_k > d+1)
\]
The two scales differ by the factor
$\frac{(\kappa_k+1)(\nu_k-d-1)}{\kappa_k(\nu_k-d+1)}$, which deviates
non-trivially from 1 when component counts $N_k$ are small.  Both approximations
improve as $N_k \to \infty$.

The NIW posterior update (\S\ref{sec:2}.2) is called at two places in the system:
the batch CAVI M-step during pretraining (\S\ref{sec:3}.2), and the
single-sample online update during play (\S\ref{sec:6}, \S\ref{sec:17}.2).  In
both cases the same four equations apply --- only the effective count $N_k$
differs (batch vs 1).

% =============================================================================
\section{Variational Bayesian Gaussian Sum (VBGS)}
\label{sec:3}
% =============================================================================

The VBGS is the agent's probabilistic ``physics engine''.  It is a
\textbf{joint Gaussian Mixture Model} trained over input-output pairs
$(s_t, s_{t+1})$: by fitting a mixture over the joint space, each component
captures a distinct dynamical regime (free flight, wall bounce, paddle bounce),
and the conditional distribution $p(s_{t+1} \mid s_t)$ is recovered analytically
by Gaussian conditioning (\S\ref{sec:4}).  This is the model used in the predict
step of the belief filter (\S\ref{sec:9}) and at step-0 of every MPPI rollout
(\S\ref{sec:15}).

Fitting is done with \textbf{Coordinate Ascent Variational Inference (CAVI)} ---
an iterative algorithm that alternates between a soft E-step and an M-step.  The
advantage over EM is that the NIW prior regularises each component automatically,
preventing degenerate solutions.

A joint GMM over $(\mathbf{x}, \mathbf{y})$ with $K$ components.
The full variational family is:
\[
q(\mathbf{Z}, \boldsymbol{\pi}, \{\mu_k, \Lambda_k\})
= q(\mathbf{Z})\,q(\boldsymbol{\pi})\prod_k q(\mu_k, \Lambda_k)
\]
with $q(\boldsymbol{\pi}) = \Dir(\boldsymbol{\alpha})$ and each
$q(\mu_k,\Lambda_k) = \cNIW(\mathbf{m}_k,\kappa_k,\nu_k,\mathbf{W}_k)$.

\subsection{CAVI E-step --- responsibilities}

For data point $\mathbf{x}_n$ and component $k$:
\[
\ln \tilde{r}_{nk}
= \psi(\alpha_k) - \psi\!\left(\sum_j \alpha_j\right)
+ \frac{1}{2}\,\E[\ln|\Lambda_k|]
- \frac{d}{2}\ln(2\pi)
- \frac{1}{2}\,\E\!\left[(\mathbf{x}_n - \mu_k)^\top\Lambda_k
  (\mathbf{x}_n - \mu_k)\right]
\]
\[
r_{nk} = \frac{\tilde{r}_{nk}}{\sum_j \tilde{r}_{nj}}
\]

\subsection{CAVI M-step --- sufficient statistics}
\[
N_k = \sum_n r_{nk}
\qquad
\bar{\mathbf{x}}_k = \frac{1}{N_k}\sum_n r_{nk}\,\mathbf{x}_n
\qquad
\mathbf{S}_k = \sum_n r_{nk}\,
  (\mathbf{x}_n - \bar{\mathbf{x}}_k)(\mathbf{x}_n - \bar{\mathbf{x}}_k)^\top
\]
\[
\alpha_k^* = \alpha_0 + N_k
\]
Then apply the NIW update from \S\ref{sec:2}.2 for each component.

\subsection{Mixing weights}

The posterior mixing proportions are:
\[
\hat\pi_k = \frac{\alpha_k}{\sum_j \alpha_j}
\]

After pretraining on $\sim$30,000 frames the ball VBGS converges to component
weights approximately $[0.29,\, 0.46,\, 0.25]$ --- wall bounces are the most
frequent regime.  The opponent VBGS converges faster (1D output) and is already
usable after a few hundred frames.  With sufficient data the M-step degenerates to
a single-sample update (\S\ref{sec:6}), enabling online refinement with negligible
computational cost.

% =============================================================================
\section{Conditional prediction}
\label{sec:4}
% =============================================================================

Training a joint GMM over $(s_t, s_{t+1})$ pairs gives a model of the
\textit{joint} distribution.  This section derives the \textit{conditional}
$p(s_{t+1} \mid s_t)$ for a single Gaussian component and then forms the mixture.
The reweighting by marginal likelihood activates the correct dynamical regime ---
the wall-bounce component gets high weight when the ball is near a wall, the
free-flight component elsewhere.

The joint NIW for component $k$ defines a joint Gaussian
$(\mathbf{x}, \mathbf{y}) \sim \cN(\boldsymbol{\mu}_k, \hat\Sigma_k)$.
Partitioned as:
\[
\boldsymbol{\mu}_k = \begin{bmatrix}\boldsymbol{\mu}_x \\ \boldsymbol{\mu}_y\end{bmatrix}
\qquad
\hat\Sigma_k = \begin{bmatrix}\Sigma_{xx} & \Sigma_{xy} \\ \Sigma_{yx} & \Sigma_{yy}\end{bmatrix}
\]
The conditional $p(\mathbf{y}\mid\mathbf{x}, k)$ is Gaussian:
\[
\mathbf{y} \mid \mathbf{x},\, k \;\sim\;
\cN\!\left(\,\boldsymbol{\mu}_y + \Sigma_{yx}\Sigma_{xx}^{-1}(\mathbf{x} - \boldsymbol{\mu}_x),\;\;
\Sigma_{yy} - \Sigma_{yx}\Sigma_{xx}^{-1}\Sigma_{xy}\,\right)
\]

\subsection{Component weights for prediction}

When conditioning on $\mathbf{x}$, component weights are re-scored using the
marginal likelihood:
\[
w_k \propto \hat\pi_k \cdot p(\mathbf{x} \mid k)
\qquad
p(\mathbf{x}\mid k) = \cN(\mathbf{x};\;\boldsymbol{\mu}_x^k,\;\Sigma_{xx}^k)
\]

\subsection{Mixture predictive mean}
\[
\hat{\mathbf{y}} = \sum_k w_k\,\boldsymbol{\mu}_{y|x}^k
\]

The mixture predictive is used in two ways downstream.  In the
\textbf{belief filter} (\S\ref{sec:9}) it provides the predicted mean and
covariance for the Kalman predict step, collapsed to a single Gaussian by
moment-matching.  In the \textbf{MPPI rollouts} (\S\ref{sec:15}) only step-0
uses the full VBGS conditional; deeper steps switch to a cheap linear model.

% =============================================================================
\section{Transition models}
\label{sec:5}
% =============================================================================

This section instantiates the VBGS architecture (\S\ref{sec:3}--\ref{sec:4}) for
the two dynamical subsystems the agent must predict: the ball and the opponent
paddle.  Each model receives weak physics-informed priors with small pseudo-counts
($\kappa_0 = 1$, $\nu_0 = d+2$) so they are overridden after a few dozen
real observations.

\subsection{Ball transition model}

Models $p(s^b_{t+1} \mid s^b_t)$ as a conditional GMM with $K=3$ components.

\textbf{Joint data:} $\mathbf{z}_t = [s^b_t;\; s^b_{t+1}] \in \R^8$,
$d_x = 4$, $d_y = 4$.

\textbf{Components and weak physics priors} (all values in $[0,1]$, $v = 0.015$):

\begin{table}[ht]
\centering
\begin{tabular}{cll}
\toprule
$k$ & Regime & Prior mean $\mathbf{m}_k$ \\
\midrule
0 & Free flight &
  $[0.5,\;0.5,\;v,\;v \;\|\; 0.5{+}v,\;0.5{+}v,\;v,\;v]$ \\
1 & Wall bounce ($\dot b_y$ flips) &
  $[0.5,\;0.1,\;v,\;v \;\|\; 0.5{+}v,\;0.1{-}v,\;v,\;{-}v]$ \\
2 & Paddle bounce ($\dot b_x$ flips) &
  $[0.85,\;0.5,\;v,\;v \;\|\; 0.85{-}v,\;0.5{+}v,\;{-}v,\;v]$ \\
\bottomrule
\end{tabular}
\end{table}

Prior hyperparameters: $\kappa_0 = 1$, $\nu_0 = d + 2 = 10$,
$\mathbf{W}_0 = 0.01\,\mathbf{I}$.

\subsection{Opponent transition model}

Models $p(o_{t+1} \mid o_t, s^b_t, a_t)$ as a conditional GMM with $K=2$
components.

\textbf{Input:}
$\mathbf{x} = [o_y,\; b_x,\; b_y,\; \dot b_x,\; \dot b_y,\; a] \in \R^6$

\textbf{Action encoding:}
\[
a = \begin{cases} 0 & \text{NOOP} \\ +1 & \text{UP} \\ -1 & \text{DOWN} \end{cases}
\]

\textbf{Joint data:} $\mathbf{z} = [\mathbf{x};\; o_{t+1}] \in \R^7$,
$d_x = 6$, $d_y = 1$.

\begin{table}[ht]
\centering
\begin{tabular}{cll}
\toprule
$k$ & Regime & Prior scale \\
\midrule
0 & Tracking (confident) & $\mathbf{W}_0 = 0.005\,\mathbf{I}$ \\
1 & Stationary / lag     & $\mathbf{W}_0 = 0.02\,\mathbf{I}$ \\
\bottomrule
\end{tabular}
\end{table}

Both transition models are pretrained offline (\code{pretrain\_models.py}) on
$\approx$30,000 frames before any live play.  The pretrained models are loaded
from \code{models/ball\_vbgs.pkl} and \code{models/opponent\_vbgs.pkl}.  The
opponent model is updated every frame regardless of the frozen/online flag.

% =============================================================================
\section{Online learning}
\label{sec:6}
% =============================================================================

After pretraining the world model captures average Pong dynamics, but individual
games can exhibit specific quirks.  Rather than freezing the world model after
pretraining, the agent performs a single CAVI E+M step on every new transition.
Because the NIW sufficient statistics accumulate data additively, this is
statistically equivalent to full batch retraining on all data seen so far ---
with zero memory overhead and negligible per-frame cost ($<0.1$ ms).

Each frame:
\begin{enumerate}
  \item Observe $(s_t, s_{t+1})$
  \item Form joint $\mathbf{z} = [s_t;\; s_{t+1}]$
  \item E-step: compute $r_{k}$ for $\mathbf{z}$
  \item M-step: update $\alpha_k$, $\mathbf{m}_k$, $\kappa_k$, $\nu_k$, $\mathbf{W}_k$
\end{enumerate}

Because the pretrained NIW statistics encode $\sim$30,000 frames, each new online
observation shifts the posterior by approximately $1/30{,}000$ --- producing slow,
noise-resistant adaptation rather than catastrophic forgetting.  A detailed
treatment of the online learning phases is given in \S\ref{sec:14}.

% =============================================================================
\section{Prior preferences}
\label{sec:7}
% =============================================================================

This section encodes \textit{what the agent wants} as a probability distribution
over observations.  In Active Inference the goal is not a hard constraint but a
soft prior $\tilde{p}(o)$: observations close to the goal are highly probable,
observations far from it are suppressed.  The Expected Free Energy (\S\ref{sec:11})
then drives the agent to minimise the KL divergence between predicted observations
and this preference distribution, producing goal-directed behaviour without an
explicit reward signal.

The agent's goals are encoded as a preferred observation distribution:
\[
\tilde{p}(o) = \cN(o \;;\; o^*,\; C)
\]
with $o = [b_x,\; b_y,\; p_y,\; o_y]$.

\subsection{Goals}

\begin{table}[ht]
\centering
\begin{tabular}{ll}
\toprule
Goal & Mechanism \\
\midrule
\textbf{G1} Ball on opponent's side &
  $o^*_{b_x} = 0.15$, tight $\sigma_{b_x} = 0.05$ \\
\textbf{G2} Player paddle tracks ball &
  $o^*_{p_y} = \mu^{b_y}$ (from belief), tight $\sigma_{p_y} = 0.03$ \\
\textbf{G3} Ball far from opponent paddle &
  $o^*_{o_y} = 0.10$, moderate $\sigma_{o_y} = 0.15$ \\
\bottomrule
\end{tabular}
\end{table}

$o^*_{b_y}$ is left loose ($\sigma = 0.20$) --- the agent does not care where
vertically the ball is, only that it reaches the opponent's side.

\subsection{Contextual target}

Goal G2 is relational: the player should track wherever the ball is.  At each step
the target is updated from the current state belief $\mu$:
\[
o^*_{p_y} \leftarrow \mu_{b_y}
\]

\subsection{Analytic expectation (used in EFE)}

When the predictive observation distribution is Gaussian
$p(o \mid q) = \cN(\mu_o, \Sigma_o)$, the expected log-preference has a closed
form:
\[
\E[\log\tilde{p}(o)]
= -\frac{1}{2}\left[
  \tr(C^{-1}\Sigma_o) + (\mu_o - o^*)^\top C^{-1}(\mu_o - o^*)
  + \log|C| + d_o\log 2\pi
  \right]
\]
This is the \textbf{instrumental value} term of the Expected Free Energy.  The
preference distribution $\tilde{p}(o)$ defined here is the fixed core;
\S\ref{sec:16} layers context-sensitivity on top by making $\sigma_{p_y}$ and
$o^*_{p_y}$ functions of the current game state.

% =============================================================================
\section{Likelihood model}
\label{sec:8}
% =============================================================================

This section specifies how the 4D RAM observation vector relates to the 6D latent
state.  The observation model $p(o\mid s)$ is a linear Gaussian (the $H$-matrix
projection), which makes the downstream Kalman belief update (\S\ref{sec:9}) exact.

\subsection{Full state vector}

The generative model operates on a 6D state that merges the ball and paddle
sub-states:
\[
s = \begin{bmatrix} b_x,\; b_y,\; \dot b_x,\; \dot b_y,\; p_y,\; o_y \end{bmatrix}^\top
\in \R^6
\]
Velocities $(\dot b_x, \dot b_y)$ are latent --- they cannot be read from a single
RAM frame.

\subsection{Observation model}

The 4D observation vector contains only the directly readable positions:
\[
o = \begin{bmatrix} b_x,\; b_y,\; p_y,\; o_y \end{bmatrix}^\top \in \R^4
\]
The likelihood is a linear Gaussian:
\[
p(o \mid s) = \cN(o \;;\; Hs,\; R)
\]
The observation matrix $H \in \R^{4 \times 6}$ selects positions from the state
(indices 0, 1, 4, 5):
\[
H = \begin{bmatrix}
1 & 0 & 0 & 0 & 0 & 0 \\
0 & 1 & 0 & 0 & 0 & 0 \\
0 & 0 & 0 & 0 & 1 & 0 \\
0 & 0 & 0 & 0 & 0 & 1
\end{bmatrix}
\]
The noise covariance $R = r\,\mathbf{I}_4$ with $r = 10^{-4}$, reflecting RAM
quantisation noise of $\pm 1$ pixel $\approx \pm 0.004$ in normalised units.

\subsection{Marginalised likelihood over a Gaussian belief}

When the agent holds a Gaussian belief $q(s) = \cN(\mu, \Sigma)$:
\[
p(o \mid q)
= \int p(o \mid s)\, q(s)\, ds
= \cN\!\left(o \;;\; H\mu,\; H\Sigma H^\top + R\right)
\]

\subsection{Kalman measurement update}

The observation $o_t$ updates the Gaussian belief via the standard Kalman
correction:
\[
S = H\Sigma H^\top + R \qquad \text{(innovation covariance)}
\]
\[
K = \Sigma H^\top S^{-1} \qquad \text{(Kalman gain)}
\]
\[
\mu^* = \mu + K\,(o - H\mu)
\qquad
\Sigma^* = (I - KH)\,\Sigma
\]
The log-likelihood of the innovation is:
\[
\log p(o) = -\frac{1}{2}\left(
  d_o \ln(2\pi) + \ln|S|
  + (o - H\mu)^\top S^{-1}(o - H\mu)\right)
\]

\subsection{Online noise estimation}

$R$ can be refined online from observed innovations
$\nu_t = o_t - H\mu_t$ via an exponential moving average:
\[
R \leftarrow (1-\eta)\,R + \eta\,\frac{1}{N}\sum_n \nu_n \nu_n^\top
\]
with learning rate $\eta = 0.01$ (Phase 4.1).

The likelihood model occupies a critical position in the information pipeline:
$H$ and $R$ appear in both the Kalman gain (\S\ref{sec:9}.2) and the EFE
epistemic term (\S\ref{sec:11}.4).  Online refinement (\S\ref{sec:8}.5) lets $R$
adapt to game-specific observation statistics over the course of a session.

% =============================================================================
\section{Belief filter (Phase 2.1)}
\label{sec:9}
% =============================================================================

The belief filter is the agent's perception system: it translates the stream of
raw RAM observations into a probability distribution
$q(s_t) = \cN(\mu_t, \Sigma_t)$ over the full 6D state, including the latent
velocities.

The filter uses a \textbf{$K=2$ mixture of Kalman filters} to handle the bimodal
predictive distribution that arises when the ball is near a wall: one component
tracks the trajectory assuming no bounce, the other assuming an imminent sign-flip
of $\dot b_y$.  This makes the filter significantly more accurate near the walls
than a single Gaussian.

The agent maintains a single Gaussian belief over the full 6D state:
\[
q(s_t) = \cN(\mu_t,\, \Sigma_t)
\]

\subsection{Predict step}

The VBGS transition models produce a mixture of Gaussians for
$p(s_{t+1} \mid s_t)$.  This is collapsed to a single Gaussian by
\textbf{moment matching}:
\[
\mu^- = \sum_k w_k\, \mu^k_{y|x}
\]
\[
\Sigma^- = \sum_k w_k\!\left(\Sigma^k_{y|x}
  + \mu^k_{y|x}{\mu^k_{y|x}}^\top\right) - \mu^-{\mu^-}^\top
\]
Ball (4D) and opponent (1D) are predicted independently; $p_y$ is
agent-controlled and carried forward unchanged.
A \textbf{process noise floor} $Q$ is added:
\[
\Sigma^-_{\text{pred}} \leftarrow \Sigma^- + Q,
\qquad
Q = \mathrm{diag}(10^{-4},\, 10^{-4},\, 10^{-4},\, 10^{-4},\, 10^{-5},\, 10^{-4})
\]

\subsection{Correct step}

The new RAM observation $o_t$ updates the belief via the Kalman measurement
equations (derived in \S\ref{sec:8}.4):
\[
S = H\Sigma^- H^\top + R
\qquad
K = \Sigma^- H^\top S^{-1}
\]
\[
\mu_t = \mu^- + K(o_t - H\mu^-)
\qquad
\Sigma_t = (I - KH)\,\Sigma^-
\]

\subsection{Why this is CAVI}

Minimising the variational free energy
$F = \E_q[\log q(s) - \log p(s,o)]$ over the Gaussian family yields exactly
these predict-correct equations.  CAVI = Kalman for linear-Gaussian models ---
no approximation is made.

\subsection{What the belief recovers}

\begin{table}[ht]
\centering
\begin{tabular}{ll}
\toprule
Quantity & Source \\
\midrule
$\mu_{b_x}, \mu_{b_y}$ &
  Smoothed ball position (less noisy than raw RAM) \\
$\mu_{\dot b_x}, \mu_{\dot b_y}$ &
  Ball velocity --- \textbf{latent}, not in any single RAM frame \\
$\Sigma$ diagonal &
  Per-dimension uncertainty --- feeds epistemic EFE term \\
\bottomrule
\end{tabular}
\end{table}

The belief $(\mu_t, \Sigma_t)$ is the single input to all downstream computation.
Accurate velocity estimation is the single most important factor in landing
prediction quality --- a 2-frame lag in detecting a bounce produces a
$\sim$0.12 normalised-unit error in the predicted landing $y$, comparable to the
paddle half-width.

% =============================================================================
\section{Validation results}
\label{sec:10}
% =============================================================================

This section reports the predictive accuracy of the pretrained VBGS transition
models before any online refinement.  After pretraining on 30,053 frames
($\approx$15 s at 2,000 fps):

\begin{table}[ht]
\centering
\begin{tabular}{lcc}
\toprule
Model & MAE & Pixels ($\times 255$) \\
\midrule
Ball position     & 0.0089 & $\approx 2.3$ px \\
Opponent position & 0.0193 & $\approx 4.9$ px \\
\bottomrule
\end{tabular}
\end{table}

Final ball component weights: $[0.29,\; 0.46,\; 0.25]$.
Component 1 (wall bounce) dominates.

The ball MAE of 2.3\,px is well below the paddle half-width of $\approx$12\,px,
giving the planner a reliable target to aim for.  The opponent MAE of 4.9\,px is
sufficient for the opponent-aware placement logic (\S\ref{sec:16}.3) to exploit.

% =============================================================================
\section{Expected Free Energy --- Phase 3}
\label{sec:11}
% =============================================================================

The Expected Free Energy (EFE) is the decision criterion that replaces the reward
function in standard reinforcement learning.  Rather than maximising a scalar
signal provided by the environment, the agent evaluates each candidate action by
asking: \textit{if I take this action, how much will the resulting observation
deviate from what I prefer, and how much uncertainty will remain about the state
of the world?}  Minimising EFE simultaneously drives the agent toward its goals
(instrumental term) and toward actions that resolve uncertainty (epistemic term).

\subsection{Full EFE decomposition}

For action $a$ the agent evaluates:
\[
G(a) =
\underbrace{-\E_{q(o|a)}[\log \tilde{p}(o)]}_{\text{instrumental}}
\;-\;
\lambda\;
\underbrace{\tfrac{1}{2}\log\frac{|S(a)|}{|R|}}_{\text{epistemic}}
\]
with $\lambda \ge 0$ controlling how much the agent values information gain.

\subsection{Predictive observation distribution}

The agent first computes a virtual one-step-ahead belief without side effects:
\[
(\mu^-_a,\;\Sigma^-_a) = \text{predict}(\mu, \Sigma, a)
\]
The predictive observation distribution is then:
\[
q(o \mid a) = \cN(o\,;\; H\mu^-_a,\; S(a)),
\qquad S(a) = H\Sigma^-_a H^\top + R
\]

\subsection{Instrumental term}

Using the analytic expectation from \S\ref{sec:7}.3 with $\Sigma_o = S(a)$ and
$\mu_o = H\mu^-_a$:
\[
G_\text{instr}(a)
= \tfrac{1}{2}\!\left[
  \tr(C^{-1} S(a))
  + (H\mu^-_a - o^*)^\top C^{-1}(H\mu^-_a - o^*)
  + \log|C| + d_o\log 2\pi
  \right]
\]
where $o^*$ is the contextual target (\S\ref{sec:7}.2) and $C$ is the preference
covariance.

\subsection{Epistemic term --- mutual information}

The epistemic term equals the mutual information $I(s'; o \mid a)$:
\[
G_\text{epist}(a)
= I(s';\,o\mid a)
= \tfrac{1}{2}\!\left(\log|S(a)| - \log|R|\right)
\]

\textbf{Derivation:}
$I(s';o) = H(o) - H(o \mid s')$.
$H(o) = \tfrac{1}{2}\log|2\pi e\, S|$ and
$H(o\mid s') = \tfrac{1}{2}\log|2\pi e\, R|$.
Their difference gives the formula above.

When $\Sigma^-_a \to 0$ (belief is certain), $S \to R$ and
$G_\text{epist} \to 0$.  Subtracting $\lambda G_\text{epist}$ from $G$ rewards
actions that resolve uncertainty.

\subsection{Action selection (1-step)}

\[
a^* = \arg\min_a\; G(a)
= \arg\min_a\;\left[G_\text{instr}(a) - \lambda\, G_\text{epist}(a)\right]
\]

$\lambda = 0$ recovers Phase~3.1 (instrumental only).
$\lambda > 0$ adds curiosity.

The 1-step EFE serves as the building block for both the horizon planner
(\S\ref{sec:12}) and the MPPI planner (\S\ref{sec:15}).  The preference target
$o^*$ and covariance $C$ entering \S\ref{sec:11}.3 are computed once per frame
and shared across all action branches.

% =============================================================================
\section{$N$-step planning horizon (Phase 3.3)}
\label{sec:12}
% =============================================================================

A single-frame EFE evaluation (\S\ref{sec:11}) is myopic: it sees only one step
ahead.  The paddle moves $\approx 0.060$ normalised units per frame; a ball
0.3 units away takes at least 5 frames to reach.  This section extends the EFE to
a full $N$-step policy tree, enabling the agent to reason that \textit{five
consecutive UPs} will close the gap while \textit{five NOOPs} will not.

\subsection{Motivation}

Single-step EFE sees only one frame ahead: the paddle moves $\pm 4$ px per step.
If the ball is 40 px away, 1-step lookahead cannot prefer \textsc{up} over
\textsc{noop} because the gap remains large regardless.  Planning $N$ steps ahead
lets the agent ``see'' that $N$ consecutive \textsc{up}s will close the gap.

\subsection{Policy tree}

A \textbf{policy} is a fixed action sequence $\pi = (a_0, a_1, \ldots, a_{N-1})$
with $a_t \in \{0, 2, 3\}$.  The number of policies is $3^N$ (27 for $N=3$, 243
for $N=5$).

\subsection{Discounted EFE for a policy}

For each policy $\pi$, the agent unrolls the belief forward $N$ steps
\textbf{in simulation} (predict steps only, no observation updates):
\[
s_0^- = (\mu_t, \Sigma_t) \quad \text{(current belief)}
\]
\[
s_{k+1}^- = \text{predict}(s_k^-, a_k), \quad k = 0, \ldots, N-1
\]
At each step $k$ the per-step EFE is evaluated:
\[
G_k(a_k) = G_\text{instr}(a_k \mid s_k^-) - \lambda\, G_\text{epist}(a_k \mid s_k^-)
\]
The policy score is the \textbf{discounted sum}:
\[
G(\pi) = \sum_{k=0}^{N-1} \gamma^k\, G_k(a_k), \qquad \gamma \in (0,1]
\]
Default: $\gamma = 0.9$.

\subsection{Action selection}

\[
\pi^* = \arg\min_\pi\; G(\pi)
\]
The agent executes only the \textbf{first action} $a^* = \pi^*_0$ and replans
at the next frame --- the ``receding horizon'' or Model Predictive Control (MPC)
strategy.

\subsection{Hybrid two-model rollout}

Naively evaluating all $3^N$ policies with full VBGS at every step costs
$\mathcal{O}(3^N \cdot N \cdot T_\text{VBGS})$, where
$T_\text{VBGS} \approx 0.7\,\text{ms}$ (677 ms per frame at $N=5$, intractable).

\textbf{Optimisation:} Step 0 uses the VBGS (wall-bounce-aware); steps
$1 \ldots N-1$ use a cheap linear-Gaussian rollout:
\[
\mu_{t+1} = F\mu_t + \delta_a,
\quad
\Sigma_{t+1} = F\Sigma_t F^\top + Q
\]
where $F$ integrates velocity and $\delta_a$ applies paddle kinematics.

Tree structure: only 3 VBGS calls (one per first action), then
$3^2 + \cdots + 3^N$ linear calls.  At $N=5$: 3 VBGS + 360 linear = \textbf{6.7
ms/frame} (100$\times$ speedup).

\subsection{Consistent preference target}

The intercept target $o^*_{p_y}$ is computed \textbf{once per frame} from the
current belief mean and shared across all branches of the tree.  Computing it
separately inside each branch introduces action-dependent noise that can corrupt
the \textsc{up} vs \textsc{down} comparison.

The tree-search horizon planner of \S\ref{sec:12} is the predecessor of the MPPI
planner (\S\ref{sec:15}), still used in \code{train\_efe.py} and
\code{validate\_efe.py} at $H=3$.

% =============================================================================
\section{Ball-intercept prior preference (Phase 3.3+)}
\label{sec:13}
% =============================================================================

This section replaces the reactive tracking target $o^*_{p_y} = \mu_{b_y}$ with
a forward-looking interception target.  By analytically simulating the ball's
trajectory from the current belief, the agent computes \textit{where the ball
will be when it arrives at the player paddle} and uses that as the preference
target from the moment the ball is served.

\subsection{Motivation}

The naive preference $o^*_{p_y} = \mu_{b_y}$ (track current ball $y$) is purely
reactive.  The agent should instead target \textbf{where the ball will arrive}
at the player's paddle, computed analytically from the current belief velocity.

\subsection{Landing prediction}

Given the current belief mean
$\mu = (b_x, b_y, v_x, v_y, p_y, o_y)^\top$:
\begin{enumerate}
  \item Clamp velocity to $|v_x|, |v_y| \le 0.06$ (suppress VBGS noise spikes
    from bounces).
  \item If $v_x \le 10^{-3}$ (ball moving away): return current $b_y$.
  \item Simulate forward: $b_x \mathrel{+}= v_x$, $b_y \mathrel{+}= v_y$ each
    step; reflect $v_y$ when $b_y \notin [0,1]$.
  \item Stop when $b_x \ge x_\text{player} \approx 0.73$ (measured contact
    point).  Return $b_y$.
\end{enumerate}

\subsection{Why this helps}

\textit{With reactive tracking} $o^*_{p_y} = \mu_{b_y}$:
the EFE only penalises the current ball--paddle gap; the ball arrives at
$b_y \ne p_y$ if it curves during flight.

\textit{With intercept targeting} $o^*_{p_y} = y_\text{land}$:
the EFE penalises the gap to the \textbf{final} position from frame 1; the paddle
starts moving immediately after serve; each new frame refines the velocity
estimate and updates the target.

\subsection{Validated results (Phase 3.3)}

\begin{table}[ht]
\centering
\begin{tabular}{lcc}
\toprule
Agent & Mean score & Best episode \\
\midrule
EFE $H=3$ (intercept) & $\mathbf{-10.8}$ & $\mathbf{-2}$ \\
EFE $H=1$ (intercept) & $-11.7$ & $-6$ \\
EFE $H=5$ (intercept) & $-14.0$ & $-8$ \\
Rule-based baseline   & $-14.8$ & $-12$ \\
\bottomrule
\end{tabular}
\end{table}

EFE $H=3$ with intercept targeting beats rule-based by \textbf{$+4$ points}
(mean).  $H=5$ underperforms because the linear rollout's constant-velocity
approximation drifts at depth 4--5 without wall-bounce corrections.

The landing prediction feeds two downstream consumers: the preference covariance
tightening in \S\ref{sec:16}.2 and the UCB sweep in \S\ref{sec:17}.4.  The raw
EMA velocity (not the Kalman $\dot b$) is used as input to avoid spike
contamination from paddle contacts.

% =============================================================================
\section{Online learning --- Phase 4}
\label{sec:14}
% =============================================================================

The pretrained models (\S\ref{sec:5}, \S\ref{sec:10}) capture average Pong
dynamics but are fixed snapshots.  This section introduces three concurrent
refinement loops that run during live play, adapting the world model to the
specific statistics of the current session.  All three operate in a streaming,
memory-free fashion by exploiting the NIW conjugate-update structure derived in
\S\ref{sec:2}.

\subsection{Online $R$ estimation (observation noise)}

The initial likelihood noise $R = 10^{-4} I$ is a fixed guess.
After each predict step, the prediction innovation is:
\[
\nu_t = o_t - H \hat{\mu}_{t|t-1}
\]
where $\hat{\mu}_{t|t-1}$ is the mixture mean after \code{predict()} but before
\code{correct()}.  Innovations are accumulated in a buffer of $B=20$ frames, then:
\[
R \leftarrow (1 - \eta_R)\,R
+ \eta_R \cdot \frac{1}{B}\sum_{i} \nu_i \nu_i^\top,
\qquad \eta_R = 0.005
\]

\subsection{Online ball VBGS refinement}

Each in-play frame produces a transition $(s_t, s_{t+1})$ approximated by the
post-correct belief mean $\mu_{t|t}[:4]$.  One CAVI E+M step is applied:
\[
r_k \leftarrow \text{E-step}(\mathbf{z}), \qquad
\theta_k \leftarrow \text{M-step}(r_k, \mathbf{z})
\]
Because the pretrained NIW statistics encode $\sim 30{,}000$ data points, each
new frame shifts the posterior by $\approx 1/30{,}000$ --- slow, noise-resistant
refinement.

\subsection{Online opponent model}

Already active since Phase~2.3: \code{Oppo\-nent\-Belief\-Tracker.update()} calls
\code{Oppo\-nent\-Transition\-VBGS.update()} every in-play frame.
No additional wiring required.

\subsubsection*{Episode structure change}

\begin{lstlisting}[caption={Frozen vs online model lifecycle across episodes.}]
Frozen (validate_efe.py)       Online (train_efe.py)
  ep 1: load -> run -> discard   ep 1: load -> run |
  ep 2: load -> run -> discard   ep 2:       run   | <- models persist
  ep 3: load -> run -> discard   ep 3:       run   |
\end{lstlisting}

\noindent\textbf{Known limitation:} the online ball VBGS path (\S\ref{sec:14}.2)
currently causes score collapse in \code{train\_efe.py} (agent scores $-20/-21$
for all episodes, misses ball on first contact).  Root cause is undiagnosed; the
ball VBGS online update is disabled by default in \code{play\_efe.py}
(\code{--learn} flag required).

% =============================================================================
\section{MPPI Receding-Horizon Planner (Phase 5)}
\label{sec:15}
% =============================================================================

The tree-search planner of \S\ref{sec:12} enumerates all $3^N$ action sequences,
which becomes intractable beyond $N \approx 5$.  This section replaces exact
enumeration with \textbf{Model Predictive Path Integral (MPPI)} control: a
stochastic sampling approach that draws $N_r = 128$ rollouts from a greedy
nominal trajectory, weights them by exponential EFE cost, and aggregates into a
soft-min estimate of the path-integral free energy.  The horizon can thereby be
extended to $H = 15$ frames while keeping per-frame compute under 7 ms.

\subsection{Motivation}

The tree-search horizon planner (\S\ref{sec:12}) grows exponentially: $3^N$
policies at horizon $N$.  At $N=15$ this is 14 million policies --- intractable.
MPPI replaces exact enumeration with a \textbf{stochastic soft-min} over $N_r$
sampled rollouts.

\subsection{Algorithm overview}

For each of the 3 root actions $a_0 \in \{0, 2, 3\}$:

\begin{enumerate}
\item \textbf{Step 0 --- VBGS root call} (wall-bounce-aware):
\[
(\hat\mu_1^{a_0},\; \hat\Sigma_1^{a_0})
= \text{VBGS-predict}(\mu_t, \Sigma_t, a_0)
\]
\[
G_0^{a_0} = G_\text{instr}(a_0) - \lambda_\text{eff}\, G_\text{epist}(a_0)
\]

\item \textbf{Greedy nominal sequence}
$\hat\pi = (\hat a_1, \ldots, \hat a_{T})$ for $T = H-1$ remaining steps:
\[
\hat a_t = \begin{cases}
\text{UP}   & \text{if } p_y^{(t)} - y^* > 0.03 \\
\text{DOWN} & \text{if } p_y^{(t)} - y^* < -0.03 \\
\text{NOOP} & \text{otherwise}
\end{cases}
\]
The deadband $\pm 0.03 = \pm\frac{1}{2} \Delta_\text{paddle}$ prevents
over-correction.

\item \textbf{$N_r$ perturbed rollouts:} each rollout $i$ copies the nominal but
  randomly replaces each step with probability $p_\text{noise} = 0.3$:
\[
\pi^{(i)}_t = \begin{cases}
\text{Uniform}\{0,2,3\} & \text{with prob } p_\text{noise} \\
\hat a_t               & \text{otherwise}
\end{cases}
\]

\item \textbf{Vectorised cost evaluation} for rollout $i$, step $t$:
\[
c^{(i)}_t = \gamma^t \left[
  \underbrace{
    \tfrac{1}{2}\!\left(
      \tr(C^{-1} S_t)
      + (\bar\mu_o^{(i,t)} - o^*)^\top C^{-1} (\bar\mu_o^{(i,t)} - o^*)
    \right)
  }_{\text{instrumental}}
  - \lambda_\text{eff}\,
  \underbrace{
    \tfrac{1}{2}(\log|S_t| - \log|R|)
  }_{\text{epistemic}}
  \right]
\]
where $\bar\mu_o^{(i,t)} = H\mu^{(i,t)}$ and $S_t$ is computed \textbf{once} per
depth (shared by all rollouts).  Total cost: $C^{(i)} = \sum_{t=1}^{T} c^{(i)}_t$.

\item \textbf{MPPI soft-min aggregation} (log-sum-exp):
\[
\tilde G^{a_0}
= -\lambda_T \log\!\left(
    \frac{1}{N_r}\sum_{i=1}^{N_r}
    \exp\!\left(-\frac{C^{(i)}}{\lambda_T}\right)
  \right)
\]

\item \textbf{Total root score} and action selection:
\[
G^{a_0} = G_0^{a_0} + \tilde G^{a_0},
\qquad
a^* = \arg\min_{a_0} G^{a_0}
\]
\end{enumerate}

\subsection{The greedy nominal is essential}

Naive MPPI uses \textsc{noop} as the nominal.  In Pong, the paddle starts far
from the target ($\sim 0.3$--$0.5$ away).  A \textsc{noop}-nominal rollout
reaches the target in 0 out of $T$ steps most of the time --- all rollouts have
similarly high cost, making the soft-min effectively uniform.  With the greedy
nominal, the nominal reaches the target in
$\lceil\text{gap}/\Delta_\text{paddle}\rceil \approx 5$--10 steps; divergent
rollouts pay extra cost; the soft-min now correctly identifies which root action
launches the optimal trajectory.

\subsection{Covariance-mean factorisation}

Because all rollouts share the same $(\hat\mu_1, \hat\Sigma_1)$ and linear
dynamics, the covariance propagation is \textbf{identical across all rollouts}.
This allows:
\begin{itemize}
  \item $S_t$, $\tr(C^{-1}S_t)$, $\log|S_t|$: computed \textbf{once per depth
    $t$}, shared by all $N_r$ rollouts.
  \item Quadratic term: vectorised over all $N_r$ rollouts as a single
    \code{einsum}.
\end{itemize}
Computational cost: $O(N_r T)$ inner products rather than
$O(N_r T \cdot d^3)$ matrix operations.

\subsection{Constants}

\begin{table}[ht]
\centering
\begin{tabular}{lll}
\toprule
Symbol & Value & Meaning \\
\midrule
$H$                  & 15    & Planning horizon (frames) \\
$N_r$                & 128   & Rollout count \\
$\gamma$             & 0.9   & Discount factor \\
$\lambda_T$          & 1.0   & MPPI temperature \\
$p_\text{noise}$     & 0.3   & Perturbation probability \\
$\Delta_\text{paddle}$ & 0.060 & Paddle step size (normalised) \\
deadband             & 0.030 & Half-step: gap below this $\to$ \textsc{noop} in nominal \\
\bottomrule
\end{tabular}
\end{table}

The MPPI planner is the outermost layer of the action-selection stack.  Its
output --- a single action index $a^* \in \{0, 2, 3\}$ --- is passed directly to
the ALE environment.  The preference target $o^*$ and covariance $C$ are prepared
by \S\ref{sec:16} from the landing prediction of \S\ref{sec:13} and the win-rate
query of \S\ref{sec:17}.

% =============================================================================
\section{Adaptive Prior Preferences (Phase 5)}
\label{sec:16}
% =============================================================================

The static preference distribution of \S\ref{sec:7} encodes the agent's
\textit{permanent} goals.  This section introduces \textit{context-sensitive}
modifications that adapt the preference to the current game state without changing
the underlying goals: the paddle-alignment precision $\sigma_{p_y}$ tightens as
the ball approaches (urgency, \S\ref{sec:16}.2); the intercept target is shifted
away from the predicted opponent position to exploit gaps in the court (placement,
\S\ref{sec:16}.3); and the target is smoothed with a bounce-aware EMA to suppress
frame-to-frame VBGS noise (\S\ref{sec:16}.5).

\subsection{Preference covariance $C$}

The agent's preferences are a diagonal Gaussian $\tilde{p}(o) = \cN(o^*, C)$:
\[
C = \operatorname{diag}(\sigma_{b_x}^2,\; \sigma_{b_y}^2,\;
                          \sigma_{p_y}^2,\; \sigma_{o_y}^2)
\]

\begin{table}[ht]
\centering
\begin{tabular}{llll}
\toprule
Dimension & Symbol & Value & Meaning \\
\midrule
$b_x$ & $\sigma_{b_x}$ & 0.05 &
  Ball must reach opponent goal (tight) \\
$b_y$ & $\sigma_{b_y}$ & 0.20 &
  Ball vertical position --- loose \\
$p_y$ (near) & $\sigma_{p_y}^\text{near}$ & \textbf{0.02} &
  Paddle must intercept ($1\sigma = \Delta_\text{paddle}/3$) \\
$p_y$ (far) & $\sigma_{p_y}^\text{far}$ & 0.04 &
  Paddle --- loose when ball distant \\
$o_y$ & $\sigma_{o_y}$ & 0.15 &
  Opponent position --- moderate \\
\bottomrule
\end{tabular}
\end{table}

\textbf{AIF principle:} $\sigma_{p_y}^\text{near} = 0.02$ encodes the agent's
\textit{goal} (hit the ball), not its motor capability.  It must never adapt to
observed errors --- that would make the agent accept misses as acceptable outcomes.

\subsection{Urgency-scaled $\sigma_{p_y}$}

$\sigma_{p_y}$ is linearly interpolated between far and near as arrival
approaches:
\[
\sigma_{p_y}(f)
= \sigma^\text{near}
+ \underbrace{\left(\frac{f - f_\text{near}}{f_\text{start} - f_\text{near}}
  \right)}_{\text{clamped to }[0,1]}
  \cdot (\sigma^\text{far} - \sigma^\text{near})
\]
where $f$ = frames to arrival, $f_\text{near} = 5$, $f_\text{start} = 50$.
The log-norm adjusts dynamically:
\[
\log|C|_\text{dyn}
= \log|C|_\text{fixed} + 2\log\sigma_{p_y} - 2\log\sigma^\text{near}
\]

\subsection{Opponent-aware placement}

When the ball is approaching ($v_x > 0$), the intercept target $y^*$ is shifted
to aim away from the predicted opponent position at ball-arrival time:
\[
y^* = y_\text{land} + \delta_\text{offset}
\]
\[
\delta_\text{offset}
= \clip\!\left(
  (o_y^\text{pred} - 0.5) \cdot 2 \cdot p_\text{max} \cdot s_\text{place},\;
  -p_\text{max},\; p_\text{max}\right)
\]
where $p_\text{max} \in [0.025, 0.07]$ is the adaptive placement magnitude and
$s_\text{place}$ fades to 0 when $f < f_\text{cutoff} = 10$ (don't sacrifice
intercept for placement).

\textbf{Intuition:} if the opponent is at $o_y = 0.3$ (above centre), the ball
is aimed high by $\delta > 0$, exploiting the gap below the opponent.

\subsection{Confidence-scaled placement}

The effective placement magnitude is scaled by ball-tracking confidence:
\[
p_\text{eff} = p_\text{max} \cdot c,
\qquad
c = \frac{1}{1 + k_c \cdot \Sigma_{b_y, b_y}}
\]
with $k_c = 30$.

\subsection{Adaptive EMA target smoothing}

Raw landing predictions fluctuate frame-to-frame from VBGS noise.  An EMA
suppresses this:
\[
y^*_\text{smooth}[t]
= \alpha \cdot y^*_\text{raw}[t] + (1-\alpha) \cdot y^*_\text{smooth}[t-1],
\quad \alpha = 0.35
\]
Time constant: $\tau = 1/(1-\alpha) \approx 3$ frames.

\textbf{Reset conditions:}
\begin{enumerate}
  \item Wall bounce detected ($v_y$ sign flip with $|dv_y| > 0.008$).
  \item Ball not approaching ($v_x \le 0$): between rallies.
\end{enumerate}

The output of \S\ref{sec:16} is the tuple $(o^*, C^{-1}, \log|C|)$ computed by
method \code{contextual\_target()} of \code{PriorPreferences}.
This is the sole interface between the preferences subsystem and the MPPI planner.

% =============================================================================
\section{Cross-Session Learning --- LongTermStats (Phase 6)}
\label{sec:17}
% =============================================================================

This section introduces persistent memory that survives between Python sessions.
The agent accumulates Bayesian sufficient statistics for two quantities: the
systematic aim bias (mean miss offset, \S\ref{sec:17}.1) and a win-rate model
indexed by the strategic context $(y^*,\; o_y^\text{arr})$ (\S\ref{sec:17}.2).
Both are serialised to disk and loaded at startup.

The win-rate model was originally a fixed-resolution 2-D Beta grid
(12~$\times$~8 cells).  Phase~4.4 replaced it with a
\textbf{continuous-feature recurrent Mixture Model (rMM)} in the AXIOM style ---
a Bayesian mixture that places NIW components and a per-component Beta(won) head
over the same 2-D feature space, growing the cluster count adaptively.

\subsection{Bayesian aim calibrator}

Tracks the systematic aim error
$\epsilon = b_y^\text{miss} - y^*_\text{pred}$ at non-fell-short misses.
Stored as \textbf{six independent EMAs} keyed by
$(\sgn(\dot b_y),\; \text{bounced} \in \{0,1\})$:
\[
\hat\mu_\epsilon^{(k)}
= \frac{\sum_{i=1}^{n_k} \epsilon_i^{(k)}}{n_k}
\quad\text{(posterior mean aim bias for key }k\text{)}
\]
\[
\text{trust}(n_k) = \min\!\left(\frac{n_k}{80},\; 0.85\right)
\]
Blended with a per-episode EMA:
\[
\text{landing\_bias}^{(k)}
\leftarrow
\text{trust}(n_k) \cdot \hat\mu_\epsilon^{(k)}
+ (1 - \text{trust}(n_k)) \cdot \text{bias}_\text{EMA}^{(k)}
\]

\subsection{Win-rate rMM}

The rMM models the joint distribution $p(\mathbf{f},\, \text{won})$ where
$\mathbf{f} = [y^*,\; o_y^\text{arr}] \in [0,1]^2$ is the strategic context at
ball-paddle contact.

\textbf{Component structure.}  Each cluster $k$ owns:
\begin{itemize}
  \item An NIW prior over the continuous features $\mathbf{f}$:
    $q(\mu_k, \Lambda_k) = \cNIW(\mathbf{m}_k,\; \kappa_k,\; \nu_k,\; \mathbf{W}_k)$
  \item A Beta(won) head:
    $q(\theta_k) = \text{Beta}(\alpha_k,\; \beta_k)$
  \item A hard-assignment count $n_k$.
\end{itemize}

\textbf{Hyperparameters.}

\begin{table}[ht]
\centering
\begin{tabular}{lll}
\toprule
Symbol & Value & Meaning \\
\midrule
$D_F$ & 2 & feature dimension $(y^*, o_y^\text{arr})$ \\
$K_\text{max}$ & 16 & maximum components \\
$\ell_\text{thresh}$ & $-2.0$ & spawn threshold on $\log p(\mathbf{f}\mid \theta_k)$ \\
$\kappa_0$ & 0.5 & NIW pseudo-count \\
$\nu_0$ & 4 & NIW dof (must be $> D_F + 1$) \\
$W_0$ & $0.0025\,\mathbf{I}$ & NIW scale ($1\sigma \approx 0.05$ in normalised units) \\
$\alpha_0,\,\beta_0$ & 2.0 & Beta prior (uniform) \\
\bottomrule
\end{tabular}
\end{table}

\textbf{Streaming update} (one observation $(\mathbf{f},\, \text{won})$).
Hard assignment, deterministic CRP-style growth:
\begin{enumerate}
  \item Compute $\ell_k = \log p(\mathbf{f}\mid \theta_k)$ for each component.
  \item \textbf{Spawn rule:} if $K = 0$, or if $K < K_\text{max}$ and
    $\max_k \ell_k < \ell_\text{thresh}$, append a new component with
    $\mathbf{m} = \mathbf{f}$, $\kappa_0, \nu_0, W_0$, $\alpha_0, \beta_0$,
    $n = 0$.
  \item Otherwise hard-assign:
    \[
    k^* = \arg\max_k\; \log\hat\pi_k + \ell_k,
    \qquad
    \hat\pi_k = \frac{n_k + 1}{\sum_j (n_j + 1)}
    \]
  \item Apply one-sample NIW M-step
    ($N_k{=}1,\, \bar{\mathbf{x}}_k = \mathbf{f},\, S_k = 0$), then update:
    \[
    \alpha_{k^*} \mathrel{+}= \1[\text{won}=1],
    \qquad
    \beta_{k^*}  \mathrel{+}= \1[\text{won}=0],
    \qquad
    n_{k^*} \mathrel{+}= 1
    \]
\end{enumerate}

\textbf{Soft responsibilities for prediction.}  Given a query $\mathbf{f}$:
\[
r_k(\mathbf{f})
= \frac{\hat\pi_k\, \exp(\ell_k)}{\sum_j \hat\pi_j\, \exp(\ell_j)}
\]

\subsection{Bayes-UCB scoring}

When the agent must choose a target $y$, it queries the rMM at
$\mathbf{f} = [y,\; \hat o_y^\text{arr}]$.

\textbf{Expected win probability:}
\[
\E[\text{won} \mid \mathbf{f}]
= \sum_k r_k(\mathbf{f})\, \frac{\alpha_k}{\alpha_k + \beta_k}
\]

\textbf{Exploration bonus (Bayes-UCB):}
\[
U(\mathbf{f})
= \E[\text{won}\mid \mathbf{f}]
+ \kappa\, \sqrt{\frac{1}{\max(n_\text{eff}(\mathbf{f}),\, 1)}}
\]
with effective observed mass
\[
n_\text{eff}(\mathbf{f})
= \sum_k r_k(\mathbf{f})\,
  \bigl((\alpha_k - \alpha_0) + (\beta_k - \beta_0)\bigr)
\]
and $\kappa = 0.3$.

\textbf{Novelty fallback.}  If $\max_k \ell_k < \ell_\text{thresh}$,
$U(\mathbf{f}) \leftarrow 0.5 + \kappa$.

\subsection{Target selection via UCB sweep}

The intercept $y$ target is chosen by Bayes-UCB argmax over 9 candidates within
the agent's reach window:
\[
y^*_{\text{cand},i}
= y_\text{land} + \frac{2(i-1) - (N-1)}{N - 1}\,\Delta_\text{reach},
\qquad i = 1, \ldots, N
\]
\[
y^* = y_{\text{cand},\, i^*},
\qquad
i^* = \arg\max_i\; U([y_{\text{cand},i},\, \hat o_y^\text{arr}])
\]
If $\Delta_\text{reach} < 10^{-3}$, the sweep is skipped and $y^* = y_\text{land}$.

\subsection{Persistence and migration}

\code{LongTermStats} is serialised inside \code{models/agent\_state.pkl} at every
episode end and on window-close in \code{play\_efe.py}.  Legacy state files with
the 2-D Beta grid are auto-migrated on load: the top $K_\text{max}$ cells by
observed mass each become a dedicated rMM cluster centred at the cell's centre;
remaining cells are replayed observation-by-observation.  A 1299-rally legacy grid
migrates cleanly in ${<}1$ s, saturating $K = K_\text{max}$.

\subsection{Episode-level placement update}

At episode end, score trend drives the placement magnitude adjustment:
\[
\text{trend} = \text{score} - \text{mean}(\text{last 10 scores})
\]
\[
\Delta p
= \alpha_p \cdot \tanh(\text{trend}/10)
  \cdot \begin{cases}
    p_\text{max} - p & \text{trend} > 0 \\
    p - p_\text{min} & \text{trend} < 0
  \end{cases}
\]
\[
p \leftarrow \clip(p + \Delta p,\; 0.025,\; 0.07),
\qquad \alpha_p = 0.05
\]

\subsection{Fell-short classification and urgency adaptation}

A miss is classified as \textbf{fell-short} when
$|p_y - y^*| > \Theta_\text{fs} = 0.090 \approx 1.5\,\Delta_\text{paddle}$.
This indicates a timing failure, not a prediction failure.

On fell-short misses, the urgency window adapts:
\[
\hat f_\text{needed}
= \frac{|p_y - y^*|}{\Delta_\text{paddle}},
\qquad
f_\text{ema} \leftarrow (1-\alpha_u) f_\text{ema} + \alpha_u \hat f_\text{needed}
\]
\[
f_\text{near} \leftarrow \text{round}(f_\text{ema}) + 2,
\qquad \alpha_u = 0.10
\]

The \code{LongTermStats} object is the only component that persists
\textit{between} Python sessions.  All other state is either frozen (preloaded
from disk) or reset per episode.

% =============================================================================
\section{Key Design Decisions and AIF Rationale}
\label{sec:18}
% =============================================================================

This section explains the non-obvious architectural choices and ties each back to
Active Inference principles.  Each subsection presents a \textit{failure mode} ---
a superficially reasonable modification that turns out to break the system ---
and explains why the current design avoids it.  Understanding these failure modes
is more valuable than simply knowing the final constants: it clarifies which
invariants the system depends on and which parameters can be safely tuned.

\subsection{Why $\sigma_{p_y}^\text{near}$ must not adapt}

After thousands of rallies with systematic aim error, \code{update\_sigma} would
have increased $\sigma_{p_y}$ to $\sim 0.05$--$0.06$.  At that scale:
\[
\Delta G(\text{NOOP vs UP})
= \frac{1}{2}\left(\frac{0.03}{0.06}\right)^2 - \frac{1}{2}\cdot 0
\approx 0.125
\]
Below VBGS noise ($\sim 0.2$--$0.5$ units) $\Rightarrow$ random action selection
$\Rightarrow$ 56\% flip rate near target.

At $\sigma = 0.02$ (fixed):
\[
\Delta G(\text{NOOP vs UP})
= \frac{1}{2}\left(\frac{0.03}{0.02}\right)^2 \approx 1.13
\]
Decisively above noise --- paddle moves correctly.

\textbf{AIF principle:} $\tilde{p}(o)$ encodes \textit{what the agent wants}
(intercept the ball).  The forward model (VBGS) should adapt to
\textit{how the world works}.  Mixing these corrupts the preference channel.

\subsection{Why the habit prior causes NOOP-lock}

At gap $= \Delta_\text{paddle}/2 = 0.03$, NOOP and UP are exactly tied in EFE.
Adding a habit bonus $-\lambda_\text{habit}$ for \code{prev\_action = NOOP}
breaks this tie unconditionally:
\[
G_\text{NOOP}^\text{eff} = G_\text{NOOP} - \lambda_\text{habit} < G_\text{UP}
\]
The paddle freezes exactly one step short of the target and misses with $\sim 50\%$
probability.  The same applies to DOWN-lock when approaching from above.

\subsection{Why the MPPI greedy nominal recovers performance}

Without a good nominal, all $N_r$ rollouts have similar high costs
$\Rightarrow$ soft-min weights are uniform
$\Rightarrow$ $\tilde G = \bar C$ (mean cost), independent of $a_0$
$\Rightarrow$ random tie-breaking at step 0.

With greedy nominal: rollouts near $\hat\pi$ converge on target; divergent
rollouts pay extra cost.  The exponential weighting in MPPI amplifies this
contrast:
\[
\text{weight}^{(i)} \propto \exp\!\left(-C^{(i)}/\lambda_T\right)
\]
Low-cost near-nominal rollouts dominate.  The soft-min correctly reflects the
value of starting with the action that enables the greedy trajectory.

These three failure modes share a common structure: each is a
\textit{plausible-sounding regularisation} that turns out to sever a critical
signal path in the Active Inference architecture.  The preference channel, the
action-selection gradient, and the nominal baseline are all load-bearing;
modifying them requires understanding what information flows through them.

% =============================================================================
\section{Roadmap}
\label{sec:19}
% =============================================================================

This table tracks which phases of the full Active Inference Pong system have been
implemented and validated.  Phases build on one another: state representation and
world models must be correct before belief filtering is meaningful, belief
filtering before EFE is correct, and EFE before planning and cross-session
learning.

\begin{table}[ht]
\centering
\begin{tabular}{clc}
\toprule
Phase & Description & Status \\
\midrule
0   & Environment, RAM observations, rule-based baseline & \yes \\
1.1 & State representation & \yes \\
1.2 & Transition models (VBGS) & \yes \\
1.3 & Likelihood model $p(o\mid s)$ & \yes \\
1.4 & Prior preferences & \yes \\
2.1 & Single Gaussian belief (CAVI / Kalman) & \yes \\
2.2 & Mixture of Gaussians belief near walls & \yes \\
2.3 & Online opponent model & \yes \\
3.1 & EFE --- instrumental value only & \yes \\
3.2 & EFE --- full with epistemic term & \yes \\
3.3 & Planning horizon $N$ steps & \yes \\
4.1 & Online $R$ estimation from innovations & \yes \\
4.2 & Online ball VBGS from live transitions & \yes \\
4.3 & Online opponent model (active since 2.3) & \yes \\
4.4 & Win-rate rMM (NIW + Beta head; replaces 2-D Beta grid) & \yes \\
5   & MPPI receding-horizon planner + greedy nominal & \yes \\
6   & Cross-session Bayesian aim calibration (LongTermStats) & \yes \\
7   & 100-episode evaluation vs rule-based and RL baseline & --- \\
\bottomrule
\end{tabular}
\caption{Implementation roadmap.}
\end{table}

% =============================================================================
\end{document}
% =============================================================================
